\documentclass[11pt,a4paper]{article}
\usepackage{amsmath,amssymb}
\usepackage{geometry}
\usepackage{hyperref}
\usepackage{parskip}
\usepackage{booktabs}

\geometry{margin=2.5cm}

\title{\textbf{\S5b --- Stratification is Forced by an Impossibility\\(draft for integration)}}

\author{Lyra}
\date{2026-09-05}

\begin{document}

\maketitle

\begin{center}
\textbf{For:} Claudius (11o1111o11oo1o1o@gmail.com)\\[4pt]
\textbf{Work-in-progress commit:} \texttt{lyra-claude/work-in-progress}
@ \texttt{1fff36d}
\end{center}

\hrule
\vspace{1em}

\section*{Stratification is forced by an impossibility, not chosen as a refinement}

\paragraph{The impossibility.}
An anytime-valid e-process $E_t$ certifies $\mathbb{E}[E_\tau] \leq 1$ at every stopping
time $\tau$.  That guarantee is \emph{marginal}---averaged over the distribution of test
points the monitor happens to see.  It does not, and cannot distribution-free, certify the
conditional statement $\mathbb{E}[E_\tau \mid \text{difficulty} = d] \leq 1$ for every
difficulty level $d$.  This is structurally analogous to the conditional-coverage impossibility
of Barber, Cand\`{e}s, Ramdas, and Tibshirani
(\href{https://arxiv.org/abs/1903.04684}{arXiv:1903.04684})%
\footnote{That result is a conformal-prediction theorem, not an e-process result; we invoke it
as a structural analogy only.}%
: by an analogous argument, with no assumptions on the data-generating distribution and in finite
samples, marginal and conditional validity cannot be had together.  The gap is structural, not an
artifact of a betting rule.
% TODO: Claudius to primary-read 1903.04684 before finalizing this cite (conformal result, analogy only).

\paragraph{Why it bites the co-failure monitor specifically.}
One is tempted to route around difficulty analytically.  The natural candidate is the detection
drift
\[
  \mu \;=\; \mathrm{KL}(P_{\mathrm{joint}} \,\|\, P_{\mathrm{marg-product}}),
\]
the Kullback--Leibler divergence between the observed joint failure law and the product of its
marginals.  Because $\mu = 0$ iff the two judges fail independently, $\mu$ is a pure dependence
functional, and it is tempting to call it ``difficulty-orthogonal by construction''---whatever
the base failure rate, $\mu$ sees only the coupling.  This is false.  Difficulty re-enters the
e-process through two channels that do not cancel---the variance of each log-likelihood increment
scales with the marginal $p(1-p)$, and estimation error in the nuisance marginals propagates into
the realized increments---so for Bernoulli failure indicators, mutual information is \emph{not} a
copula invariant, and difficulty cannot be conditioned out for free.  The illustration is direct:
fix the odds ratio (the coupling) and sweep the marginal failure rate from $0.05$ to $0.5$,
and $\mu$ moves by an order of magnitude.  The only sense in which $\mu$ ``conditions out'' difficulty
is circular: define ``same coupling'' to mean ``same $\mu$'', which defines the problem away rather
than solving it.

\paragraph{Consequence.}
So the drift law, $\text{latency} \propto 1/\mu$, survives as the quantitative heart of the
monitor---but not as a difficulty-robust one.  Its non-robustness \emph{is} the
conditional-coverage gap: no distribution-free functional of the stream can deliver per-difficulty
validity for free.  Stratification is what the impossibility leaves open---the one escape it
does not close.\footnote{The direct escape aside, shape constraints and DKW corrections may offer
alternatives we have not yet closed.}  We stratify not to refine the estimator but because the theorem forbids the
alternative.

\paragraph{The positive construction.}
Partition the difficulty axis into $K$ strata and run a separate betting e-variable within
each.  On items in stratum $k$,
\[
  e \;=\; 1 + \lambda\bigl(U - V - \delta_k\bigr),
\]
where $U = \mathbf{1}\{\text{both judges fail on item }i\}$ and $V = W_i^s W_j^t$ pairs the two
judges across \emph{different} items $s \neq t$ drawn from items other than the current item,
$s, t \neq i$, so the marginals multiply exactly,
$\mathbb{E}[V] = ab$, with no fitted nuisance.  The slack $\delta_k = 2\varepsilon$ absorbs the
worst-case two-sided excursion of the stratum's base rate, and the betting parameter is confined
to $\lambda \in [0, 1/(1+2\varepsilon)]$ so the e-variable stays non-negative.  The product
over strata (or a mixture across $k$) is again a valid e-process by Ville.  Conditional
coverage is recovered by construction---one certificate per stratum---at the $O(K)$ cost the
impossibility theorem says such coverage must carry.

\paragraph{The practical difference.}
On a synthetic panel where difficulty and coupling are confounded, the three methods behave
as follows:

\begin{center}
\begin{tabular}{llc}
\toprule
Method & Description & Type-I error \\
\midrule
A & Stratified cross-item monitor & $0.00$ \\
B & Unstratified cross-item monitor & $0.23$--$0.44$ \\
C & Na\"{\i}ve plug-in (fitted marginals) & $0.59$--$0.70$ \\
\bottomrule
\end{tabular}
\end{center}

Stratification is not tightening a bound; it is the difference between a valid test and an
invalid one.  (Synthetic; a proof of necessity, not an empirical effect size.)

\paragraph{A second conditioning variable, and a distinction we keep.}
Difficulty is not the only axis on which a panel's ballots carry structure.  Shu's pivotality
analysis (\href{https://arxiv.org/abs/2608.06940}{arXiv:2608.06940}) conditions on the vote
margin $m_i = |2s_i - k|$: nearly all of a panel's recoverable accuracy sits on the $\approx
16\%$~\footnote{The $\approx 16\%$ figure is Shu's; pending our own primary read (arXiv:2608.06940).} of queries where a single flip would change the verdict, and, unlike difficulty, margin
needs no labels to compute.  It is tempting to treat difficulty and margin as interchangeable
stratification axes governed by one impossibility.  They are governed by the same
impossibility---neither can be conditioned out for free---but they are not interchangeable, and
the difference is where our construction lives.

Difficulty is \emph{ancillary} to the tested dependence: conditioning on it leaves the two
judges' ballots conditionally independent, so $\mathbb{E}[V] = ab$ is preserved stratum by
stratum.  Margin is \emph{not} ancillary; it is near-sufficient for the co-failure alternative,
because co-failure \emph{is} excess agreement and margin is a function of the agreement count.
Conditioning on the ballot sum $s_i = c$ pins each judge's marginal to $c/k$ and induces a
negative within-item covariance,
\[
  \mathrm{Cov}(W_i, W_j \mid s_i = c)
  \;=\; -\frac{(c/k)(1 - c/k)}{k-1} \;<\; 0,
\]
which breaks the exact factorization $\mathbb{E}[V] = ab$ that the whole margins-free
construction rests on---and drains the test's power precisely at unanimity, where co-failure
is strongest.  The distinction is that difficulty is \emph{exogenous} to the ballots---so we
stratify on it---whereas margin is \emph{endogenous}, a deterministic function of the very
agreement count the test interrogates.  We therefore keep margin strictly as a \emph{post-hoc
validity diagnostic}---a label-free check that Type-I control does not depend on where on the
margin scale the panel sits---and \emph{not} as a stratification variable co-equal with
difficulty.  This resolution is ours.  Item independence closes the remaining step: since $V = W_i^s W_j^t$
pairs judges across items $s, t \neq i$, and items are independent by construction, the ballot
sum $s_i = c$ is conditionally uninformative about $W^s$ and $W^t$; the joint distribution of
$V$ is unchanged by the conditioning, and the covariance structure of $U$ and $V$ decouples
exactly as required.

\paragraph{The perspectivist objection scopes the null, it does not refute it.}
A perspectivist objection---in the tradition of subjectivity-aware NLP annotation (Plank;
Xu \& Jurgens)---holds that on genuinely subjective or ambiguous items, disagreement among
judges is not error but legitimate variation in perspective, so a low effective judge count on
such items is no monitoring failure.  We grant this, and it sharpens the null's scope rather than
threatening the monitor's validity.  The co-failure monitor targets items with a \emph{determinate} ground truth,
where correlated judge failure is a genuine reliability defect; on genuinely-indeterminate items
an effective count of $\approx 2$ is not a pathology to flag but real perspective spread, and the
monitor should stay silent.  The remedy is exactly the move we made for difficulty---stratify by
ground-truth-determinability---because determinability is \emph{exogenous} to the ballots (a
property of the item, like difficulty), hence a legitimate stratification variable.  We concede that
determinability, like difficulty, is not directly observed and must be supplied by external
(meta-)annotation, which places it on the same exogenous footing rather than making it freely
readable off the stream.  Conditioning
on it separates co-failure on determinate items, which we monitor, from legitimate perspective
diversity on indeterminate items, which is not our target.  The perspectivist critique thus scopes
the null; it does not refute it.

\paragraph{Alternatives under investigation (brief).}
Two alternatives to explicit stratification are under investigation and not yet adopted:
shape-constrained e-processes (Wang--Ramdas, \href{https://arxiv.org/abs/2604.20612}{arXiv:2604.20612}),
which would monitor whether the failure-rate curve keeps its monotone/unimodal shape in
difficulty rather than stratifying it out; and DKW-corrected conditional test martingales
(\href{https://arxiv.org/abs/2602.13848}{arXiv:2602.13848}).  Both are flagged pending primary
reading and are mentioned only for completeness.

\end{document}
